\documentclass[11pt,letterpaper]{article}
\PassOptionsToPackage{hyphens}{url}
\usepackage[margin=.83in,top=.8in,bottom=.8in,headheight=15pt]{geometry}
\usepackage[T1]{fontenc}
\usepackage{lmodern,amsmath,amssymb,mathtools,bm,booktabs,array,graphicx,tikz}
\usetikzlibrary{arrows.meta,positioning,calc,decorations.pathreplacing,patterns}
\definecolor{ink}{HTML}{17344D}\definecolor{accent}{HTML}{087E8B}\definecolor{pale}{HTML}{EDF5F5}\definecolor{gold}{HTML}{AD6B19}
\providecommand{\dd}{\,\mathrm{d}}\providecommand{\R}{\mathbb{R}}
\providecommand{\E}{\mathbb{E}}\providecommand{\Var}{\operatorname{Var}}
\providecommand{\Cov}{\operatorname{Cov}}\providecommand{\diag}{\operatorname{diag}}
\providecommand{\norm}[1]{\left\lVert#1\right\rVert}
\providecommand{\abs}[1]{\left|#1\right|}
\usepackage{microtype,hyperref,fancyhdr,enumitem,amsthm,sectsty,needspace}
\allsectionsfont{\raggedright}
\hypersetup{colorlinks=true,linkcolor=ink,urlcolor=accent,pdfauthor={Zan Ahmad},pdftitle={Functions, array axes, and planar trajectories / Review}}
\setlength{\parindent}{0pt}\setlength{\parskip}{6pt}\linespread{1.07}
\setlist{itemsep=3pt,topsep=4pt,leftmargin=*}
\emergencystretch=2em\clubpenalty=10000\widowpenalty=10000
\pagestyle{fancy}\fancyhf{}\fancyhead[L]{\small Scientific Computing in Python / 2027}\fancyhead[R]{\small Lecture 2}
\fancyfoot[L]{\footnotesize Zan Ahmad}\fancyfoot[R]{\footnotesize\thepage}\renewcommand{\headrulewidth}{.3pt}
\newtheorem{proposition}{Proposition}\newtheorem{theorem}{Theorem}\newtheorem{lemma}{Lemma}
\theoremstyle{definition}\newtheorem{definition}{Definition}\newtheorem{example}{Example}
\numberwithin{equation}{section}
\begin{document}
\section*{Functions, array axes, and planar trajectories --- review sheet}
Use this after attempting the worked examples and practice. Explain each claim in words before calculating.

\section*{Key ideas}
\begin{itemize}
\item A function exposes an input--output contract.
\item Return values can feed later calculations.
\item Loops make repeated index meanings visible.
\item Name every axis before broadcasting or reducing.
\item A rotation uses matrix multiplication and preserves lengths.
\item A threshold is a declared decision rule, not a scientific truth.
\end{itemize}
\section*{Common errors and repairs}
\textbf{Treat the final array axis as time automatically.} Here the final axis is coordinates; time is axis 1.

\textbf{Divide a batch by a duration array with no axis reasoning.} Align durations as (1, intervals, 1).

\textbf{Call every norm an observed path speed.} It is the magnitude of average velocity over sampled intervals.

\section*{Retrieve without your notes}
\begin{itemize}
\item What shape follows time differencing of (8,10,2)?
\item Which of > and >=\allowbreak{} includes equality?
\item What remains invariant under a rigid rotation?
\end{itemize}
\section*{Readiness check}
\begin{itemize}
\item Write and test a function with inputs, a return value and a declared invalid-input case.
\item Trace a loop and a strict comparison at its boundary.
\item Label trial, time and coordinate axes in a three-dimensional NumPy array.
\item Compute interval-average velocity components and their magnitudes with explicit broadcasting.
\item Use a matrix rotation and check that it preserves displacement lengths.
\end{itemize}
\end{document}
